\section{Discussion}
\label{sec:discussion}

The relaxation and the entering construction serve different parts of the search.
The rate relaxation allocates duties across robot types and sites, while the construction brings a fresh fleet into that rotation.
Removing the construction raises the gap by $2.3$--$3.2$ percentage points, and removing the search raises it by $1.1$--$2.7$ points (Table~\ref{tab:ablation}).
Thus, the construction supplies a starting schedule, and the deviations reduce the resources needed to run it.
The search changes both the fleet and the decisions that determine its initial phase distribution.

The deviation window and the evaluation horizon have different roles.
On the $40$ instances used for the window experiment, extending the window from $0.5S_{\max}$ to $2S_{\max}$ raises the mean gap from $12.4\%$ to $13.2\%$ under the same time budget.
At a fixed deviation rate, the longer window adds genes and spreads the search over more decisions.
The shorter window concentrates those decisions near the fresh start.
Their effects continue after the genes stop acting because they change which robots are ready and where they are in their maintenance cycles.
Evaluation must follow those effects: $157$ of the $240$ schedules selected using only the first $720$ hours run short later.
A short encoding therefore works together with a long-run check.

The remaining gap combines the relaxation of the resource counts and trajectories with the quality of the schedules found.
The start-up LP allows fractional paths and stops coverage at $H$; C-RDD assigns whole robots and tools and follows the resulting schedule beyond that horizon.
Fig.~\ref{fig:bracket} shows that the rotation and the entering construction use similar numbers of resources, so much of the gap is already present before the fresh fleet enters.
This points to the choice of rotation as a place for further improvement.
Searching over several duty allocations from the rate relaxation could change both the steady fleet and the cost of entering it.
On the lower side, linking the terminal state of the start-up relaxation to a continuing rotation could account for coverage after $H$.
